\chapter{Ensemble}
\label{chap:ensemble}

Many models already exist, fully trained, each producing an out-of-fold probability for every
window. Rather than train anything new, this chapter asks a cheaper question: does simply
averaging the probabilities of two or more already-trained models do better than the best one
alone?

\section{Method}

Each ensemble blends the out-of-fold probabilities of two or more already-trained models (and,
in several rows below, one CNN) by averaging them --- no retraining involved. Two different
threshold conventions are reported, because they answer different questions:

\begin{itemize}[leftmargin=1.4em]
  \item \textbf{Held-out-fold} (the honest number): for each of the 5 outer folds, the detection
    threshold is swept using only the \emph{other} four folds' predictions, never the fold being
    scored.
  \item \textbf{Optimistic}: a single threshold swept across all out-of-fold predictions at
    once --- the same convention used for the ``best threshold'' column in
    \Cref{tab:classical-results}, kept here only for comparison.
\end{itemize}

\section{Results}

\begin{table}[H]
  \centering
  \small
  \caption{Ensemble blends, held-out-fold threshold, against the two single-model baselines.}
  \label{tab:ensemble-results}
  \begin{tabular}{p{6.4cm}rrrr}
    \toprule
    Ensemble & Event F1 & Found/93 & FA & SzCORE F1 \\
    \midrule
    SVM, channel dropout (single model, new best) & \textbf{0.901} & 88 & 18 & --- \\
    SVM Unbalanced (single model, original baseline) & 0.894 & 86 & 16 & --- \\
    SVM Unbalanced + ShallowConvNet Augmented & 0.901 & 90 & 19 & 0.900 \\
    SVM, channel dropout + ShallowConvNet, amplitude & 0.896 & 89 & 19 & 0.892 \\
    SVM Unbalanced + Random Forest Unbalanced & 0.894 & 85 & 15 & 0.894 \\
    SVM, channel dropout + Random Forest Unbalanced & 0.823 & 84 & 31 & 0.866 \\
    SVM, channel dropout + Random Forest, channel dropout & 0.817 & 81 & 29 & 0.846 \\
    SVM Unb.\ + RF Unb.\ + KNN Unb. & 0.874 & 83 & 18 & 0.881 \\
    SVM Unbalanced + KNN Unbalanced & 0.871 & 81 & 15 & 0.872 \\
    2$\times$SVM Unb.\ + Random Forest Unb. & 0.849 & 83 & 23 & 0.870 \\
    \bottomrule
  \end{tabular}
\end{table}

\textbf{No ensemble beats the best single model.} SVM with channel-dropout augmentation
(\Cref{chap:augmentation}) reaches event F1 0.901 on its own. The best ensemble found, SVM
Unbalanced blended with the augmented ShallowConvNet CNN, reaches exactly the same score, 0.901
--- a genuine tie, not an improvement, and it does so with a different, more expensive recipe (a
second model to train and maintain) for the same result a single classical model already gives.

Blending the \emph{new} best SVM with anything else actually makes things worse, not better:
paired with the best CNN it drops to 0.896; paired with Random Forest (either version) it drops
further, to 0.817--0.823. This makes sense in hindsight --- the earlier ensemble gains came from
combining models that made \emph{different kinds} of mistakes (SVM unbalanced is conservative
and precise, the CNN and Random Forest catch different seizures). The new best SVM already
captures much of that diversity internally, through the channel-dropout training itself, so
blending it with another model adds noise more than it adds new information.

The older SVM+Random Forest blend (both unbalanced) still ties the original 0.894 baseline, with
a different mix of hits and misses (85 found, 15 false alarms, against the SVM alone's 86 and
16) --- a result worth remembering but not a reason to prefer it over either single-model result
above.

\begin{summarybox}
Blending already-trained models, with no retraining, was tried extensively. The best blend found
(SVM unbalanced + an augmented CNN) ties the project's best single model (0.901) but never beats
it, and blending the new best SVM (channel-dropout augmentation) with anything else makes results
worse. The practical conclusion: a single, well-chosen model --- SVM with channel-dropout
augmentation --- is both the simplest and the best-performing option in this project, with no
ensemble complexity needed to reach the same score.
\end{summarybox}
